\section{A zero diagonal annihilates bounded hypercohomology}

Let $\mathcal A$ be an abelian category and let $K^\bullet$ be a
cochain complex with terms zero outside a finite interval.  For an object
$A$ of $\mathcal A$, write
\[
 \mathbb H_A^t(K^\bullet)
 =\operatorname{Hom}_{D(\mathcal A)}(A[-t],K^\bullet).
\]
For abelian sheaves and the constant free abelian sheaf $A=\mathbb Z$,
this is sheaf hypercohomology.  A sheaf $B$ in degree $i$ contributes
$\operatorname{Ext}^{t-i}(A,B)$ when $t-i\geq0$.

\begin{lemma}[The exact sequence of consecutive brutal tails]
\label{lem:kn-hypercohomology-brutal-tail}
\lean{EnvelopeProperty.Vanishing.HypercohomologyBrutalTail.sequence_shortExact}
\leanok
For an integer $a$, let $K_{\geq a}^\bullet$ have the original terms
and differentials in degrees at least $a$ and zero terms below $a$.
There is a degreewise short exact sequence
\[
 0\longrightarrow K_{\geq a+1}^\bullet
 \longrightarrow K_{\geq a}^\bullet
 \longrightarrow K^a[-a]\longrightarrow0.
\]
\end{lemma}
\begin{proof}
The first arrow is the identity in every degree above $a$ and zero
otherwise.  The second arrow is the identity in degree $a$ and zero
otherwise.  Both commute with the original differentials: the first
has zero source below $a+1$, while the second has zero target outside
degree $a$.  In each degree below $a$ all terms are zero; in degree
$a$ the sequence is $0\to0\to K^a\xrightarrow{1}K^a\to0$; in each
degree above $a$ it is $0\to K^i\xrightarrow{1}K^i\to0\to0$.
These are short exact, proving the assertion in the category of
complexes.
\end{proof}

\begin{theorem}[Bounded zero-diagonal hypercohomology]
\label{thm:kn-hypercohomology-zero-diagonal}
\lean{EnvelopeProperty.Vanishing.HypercohomologyDiagonal.subsingleton_hom_of_bounded_terms,EnvelopeProperty.Vanishing.HypercohomologyDiagonal.subsingleton_single_hypercohomology_of_H}
\leanok
\uses{lem:kn-hypercohomology-brutal-tail}
Suppose $K^i=0$ outside $a\leq i<b$.  If
\[
 \operatorname{Hom}_{D(\mathcal A)}(S,K^i[-i])=0
 \quad(a\leq i<b),
\]
then $\operatorname{Hom}_{D(\mathcal A)}(S,K^\bullet)=0$.
In particular, for a bounded complex of abelian sheaves and an integer
$t\geq b-1$, vanishing of every $H^{t-i}(K^i)$ on this whole diagonal
implies $\mathbb H^t(K^\bullet)=0$.
\end{theorem}
\begin{proof}
The brutal tail starting in degree $b$ is zero.  Descend through the
finitely many integers $b-1,\ldots,a$.  At each step apply
$\operatorname{Hom}(S,-)$ to the distinguished triangle associated
to the short exact sequence of consecutive tails.  Exactness gives
\[
 \operatorname{Hom}(S,K_{\geq i+1}^\bullet)
 \longrightarrow\operatorname{Hom}(S,K_{\geq i}^\bullet)
 \longrightarrow\operatorname{Hom}(S,K^i[-i]).
\]
The outer groups are zero by induction and hypothesis, hence so is
the middle group.  The final tail is canonically the original complex,
since all its terms below $a$ are zero.  For sheaves take
$S=\mathbb Z[-t]$.  Shifting both arguments identifies the last
single-term Hom group with
$\operatorname{Ext}^{t-i}(\mathbb Z,K^i)=H^{t-i}(K^i)$.
No finite-dimensionality, duality, or decomposition assumption enters
this argument.
\end{proof}
